\section{A necessary depth for bounded absolute rate}
\label{sec:critical-entropy-transport}

The common limiting rate does not mean that a bounded rate is attained
at every depth. Gaussian entropy transport gives a complementary
constraint. This is a statement about representation rates, not a
sample-complexity lower bound.

For $y=(y_0,\ldots,y_d)\in\{-1,1\}^{d+1}$ define the odd extension
\[
 \widetilde f(y)=y_0 f(y_1/y_0,\ldots,y_d/y_0),
\]
and let $k=k_{\rm aug}(f)$ be its number of relevant coordinates.
Here a coordinate is relevant when some edge parallel to that coordinate
has opposite endpoint labels. Since $\widetilde f$ is odd, $1\le k\le d+1$.
The count is unchanged by permutations of the augmented coordinates.
In particular, constants and dictators both have $k=1$. A parity of
original degree $j$ has $k=j$ when $j$ is odd and $k=j+1$ when $j$ is even.

\begin{theorem}[Entropy transport and depth]
\label{thm:critical-entropy-transport}
Let $J=J_L(f)$ and $h=L-2$. Put
\[
 B_h(J)=
 \begin{cases}(1+\sqrt{2J/h})^{2h},&h\ge1,\\1,&h=0.\end{cases}
\]
Then
\begin{align}
 k_{\rm aug}(f)&\le64J B_h(J),
 \label{eq:critical-transport-relevance}\\
 d&\le8J\bigl(1+\sqrt{2J/k_{\rm aug}(f)}\bigr)^2 B_h(J).
 \label{eq:critical-transport-dimension}
\end{align}
In particular $B_h(J)\le\exp(2\sqrt{2hJ})$. If $J_L(f)\le K$ and
$d>8K(1+\sqrt{2K})^2$, then
\[
 L\ge2+\frac{1}{8K}
 \left[\log\frac{d}{8K(1+\sqrt{2K})^2}\right]^2.
\]
Thus every target, including a constant, requires depth
$\Omega((\log d)^2)$ to have a uniformly bounded critical rate as
$d\to\infty$.
\end{theorem}

We use the Gaussian transport inequality in its correlated form. If
$\eta\ll\mathcal N(0,Q)$ has entropy $D$, there is a coupling $(Z,G)$
with these two marginals such that
\begin{equation}
 \mathbb E\|Q^{+1/2}(Z-G)\|^2\le2D.
 \label{eq:critical-gaussian-transport}
\end{equation}
Both variables lie in $\operatorname{range}Q$. This is the standard
Gaussian quadratic transport inequality~\cite{talagrand}, applied after whitening on that
range; singularity causes no change. In particular one common coupling
works for all input coordinates, rather than separate coordinatewise
couplings.

\begin{proof}[Proof of Theorem~\ref{thm:critical-entropy-transport}]
Take any admissible variational path of total cost
$A=r+\sum_{s=1}^{L-1}D_s<\infty$, where $r$ is its readout cost.
Extend every feature oddly to the augmented cube as above. This is an
injective linear transformation of the feature vector and preserves
relative entropy. The Gaussian reference and the tanh map commute with
this extension. Finite entropy ensures that every feature coordinate is
in $(-1,1)$, so coordinatewise $\operatorname{atanh}$ lifts each controlled
feature law to its preactivation law with exactly the same entropy.

For a later hidden layer with input second moment $Q$, apply
\eqref{eq:critical-gaussian-transport}. For any two augmented inputs
$y,y'$, put $b=e_y-e_{y'}$. Cauchy--Schwarz on
$\operatorname{range}Q$ and Minkowski give
\[
 \|b^TZ\|_{L^2}
 \le\|b^TG\|_{L^2}+\|b^T(Z-G)\|_{L^2}
 \le(1+\sqrt{2D_s})\sqrt{b^TQb}.
\]
Since tanh is $1$-Lipschitz, the $L^2$ distance of the two output
features is at most this bound. Hence each later layer multiplies
pairwise feature distances by at most $1+\sqrt{2D_s}$. Write
\[
 P=\prod_{s=2}^{L-1}(1+\sqrt{2D_s}),
\]
with the empty product equal to one.

At the first layer, the field $y\mapsto W\cdot y$ identifies $W$
uniquely. Its controlled law has entropy $D_1$ relative to
$\mathcal N(0,I_{d+1}/d)$. The Gaussian transport coupling therefore
satisfies $\mathbb E\|W-G\|^2\le2D_1/d$. If $I$ is the set of the
$k$ relevant augmented coordinates, projection and Minkowski give
\begin{equation}
 \sum_{i\in I}\mathbb EW_i^2
 \le\frac1d\bigl(\sqrt k+\sqrt{2D_1}\bigr)^2,
 \qquad
 \mathbb E\|W\|^2
 \le\frac1d\bigl(\sqrt{d+1}+\sqrt{2D_1}\bigr)^2.
 \label{eq:critical-transport-first-trace}
\end{equation}
A first-layer edge in direction $i$ has feature distance at most
$2\sqrt{\mathbb EW_i^2}$.

The readout admits a representation $z(y)=\mathbb E[a(H)H(y)]$
with $\mathbb Ea(H)^2=2r$: one may take
$a(H)=z^TQ^+H$. Each relevant direction has a bichromatic edge, and
the output margin makes the output difference on that edge at least
one. Consequently
\[
 1\le2\sqrt{2r}\,P\sqrt{\mathbb EW_i^2}\qquad(i\in I).
\]
Summing the squared inequalities and applying the first estimate in
\eqref{eq:critical-transport-first-trace} proves the sharper budget form
\begin{equation}
 d\le8rP^2\bigl(1+\sqrt{2D_1/k}\bigr)^2.
 \label{eq:critical-transport-exact-budget}
\end{equation}
Using the full trace instead gives
\[
 k\le\frac{8rP^2}{d}
       \bigl(\sqrt{d+1}+\sqrt{2D_1}\bigr)^2.
\]
If $A<d$, its right side is at most $64AP^2$. If $A\ge d$,
the same conclusion follows from $k\le d+1\le2A$.

The function $u\mapsto\log(1+\sqrt{2u})$ is concave on
$[0,\infty)$. For $h\ge1$, Jensen's inequality and
$\sum_{s=2}^{L-1}D_s\le A$ give
\[
 P^2\le(1+\sqrt{2A/h})^{2h}
       \le e^{2\sqrt{2hA}}.
\]
The same assertion with $P=1$ holds for $h=0$. Substituting
$r,D_1\le A$ in the preceding bounds and taking $A\downarrow J$
proves both inequalities. The final depth bound follows by taking
logarithms of \eqref{eq:critical-transport-dimension} and using $k\ge1$.
\end{proof}

\begin{remark}[What the bound does and does not control]
For fixed rate budget $K$, \eqref{eq:critical-transport-relevance}
also requires depth $\Omega((\log k)^2)$ when $k\to\infty$.
Any Boolean predictor that has uniform-distribution error strictly less
than $1/2$ on a parity of $j$ specified coordinates must depend on each
of those coordinates. Thus the same relevant-coordinate restriction
applies to every such predictor. Neither statement compares posterior
masses of functions consistent with a training sample. A bound on
pairwise feature distances also does not by itself bound the number of
possible correlated feature laws, so it cannot be read as a learning
theorem. The ambient-dimension restriction above is compatible with
Theorem~\ref{thm:critical-rate-flattening}, whose sufficient depth is
polynomial in $d$.
\end{remark}
